% ID: SEL-D0037-OA-E03
% ESTADO: cotejada. Segunda pasada visual y borrador compilado revisados.
% Fuente: Selectividad/matematicas_aplicadas/Ua2MATEMÁTICAS APLICADAS examen.pdf
% SHA-256: a56739d297d4dc220d6898b24cf729dc556401a6bb292b8ca7cfb58c01c31651
% Página: 1; opción A; ejercicio original 3.

El 70\% de los clientes de un supermercado realizan las compras en el local y el resto de los clientes las realizan por internet. De las compras realizadas en el local, sólo el 30\% supera los 100 \texteuro{}, mientras que de las realizadas por internet el 80\% supera esa cantidad.
\begin{enumerate}[label=\alph*)]
\item \textbf{(1.5 puntos)} Elegida una compra al azar, \textquestiondown{}cuál es la probabilidad de que supere los 100 \texteuro{}?
\item \textbf{(1 punto)} Si se sabe que una compra supera los 100 \texteuro{}, \textquestiondown{}cuál es la probabilidad de que se haya hecho en el local?
\end{enumerate}
